\begin{center}
  {\Large \bfseries TÓM TẮT NỘI DUNG}
\end{center}
\vspace{0.5cm}

Việc khai thác và sử dụng bền vững các nguồn tài nguyên sinh học biển đóng vai trò đặc biệt quan trọng trong việc bảo tồn đa dạng sinh học và phát triển kinh tế lâu dài, bởi việc khai thác quá mức trong thời gian dài có thể dẫn đến sự cạn kiệt trữ lượng sinh khối và đe dọa trực tiếp sự tồn vong của toàn bộ hệ sinh thái. Vịnh Chesapeake (Hoa Kỳ) là một trong những hệ sinh thái ven biển phức tạp và năng động nhất thế giới, nơi mà việc duy trì nghề cá bền vững là nhân tố then chốt cho cả sự thịnh vượng kinh tế lẫn an ninh sinh thái của toàn vùng duyên hải.

Báo cáo này nghiên cứu và tái hiện toàn diện bài báo khoa học \textit{``Bioeconomic analysis of harvesting within a predator--prey system: A case study in the Chesapeake Bay fisheries''} của Iordanka N. Panayotova, John Herrmann và Nathan Kolling (công bố trên tạp chí \textit{Ecological Modelling}, 2023). Nghiên cứu xây dựng mô hình toán học dựa trên hệ sinh thái để phân tích động lực học quần thể có khai thác của hai loài cá trọng yếu trong vịnh Chesapeake: cá mòi dầu Đại Tây Dương (\textit{Brevoortia tyrannus}) đóng vai trò là con mồi (\textit{prey}) và cá vược sọc (\textit{Morone saxatilis}) đóng vai trò là động vật săn mồi (\textit{predator}).

Mô hình sử dụng hệ phương trình vi phân Lotka--Volterra tổng quát phi tuyến, được hiệu chỉnh và khớp tham số với chuỗi dữ liệu thực tế về độ phong phú cá thể từ năm 1995 đến năm 2017. Báo cáo phân tích tính đặt đúng về mặt sinh học của mô hình (tính dương và tính bị chặn của nghiệm trong miền bất biến $\Omega$), thiết lập các điều kiện tồn tại và ổn định của các điểm cân bằng kinh tế sinh thái (bionomic equilibria). Đặc biệt, nghiên cứu đi sâu so sánh ba chiến lược quản lý khai thác chủ chốt:
\begin{enumerate}
    \item \textbf{Sản lượng bền vững tối đa (MSY - Maximum Sustainable Yield):} Tối đa hóa sản lượng sinh khối bắt được trong điều kiện duy trì trữ lượng quần thể.
    \item \textbf{Sản lượng kinh tế tối đa (MEY - Maximum Economic Yield):} Tối đa hóa lợi nhuận ròng (tô kinh tế) từ việc khai thác sau khi trừ chi phí.
    \item \textbf{Sản lượng tối đa hóa độ đàn hồi sinh thái (RMY - Resilience Maximizing Yield):} Chiến lược quản lý ưu tiên khả năng tự phục hồi nhanh nhất của hệ thống về trạng thái cân bằng sau những nhiễu loạn môi trường.
\end{enumerate}

Cuối cùng, phương pháp vi phân trực tiếp được áp dụng để phân tích độ nhạy cục bộ của các tham số động thái, qua đó xác định các yếu tố có tác động chi phối mạnh nhất đến trữ lượng hai loài cá. Các kết quả này cung cấp luận cứ toán học định lượng vững chắc giúp các nhà hoạch định chính sách cân bằng hài hòa giữa lợi ích kinh tế trước mắt và mục tiêu bảo tồn hệ sinh thái bền vững trong dài hạn.

\vspace{0.4cm}
\noindent \textbf{Từ khóa:} Mô hình hóa toán học, Phương trình vi phân phi tuyến, Hệ động lực con mồi -- kẻ săn mồi, Vịnh Chesapeake, Kinh tế sinh thái (Bioeconomics), Sản lượng bền vững tối đa (MSY), Sản lượng kinh tế tối đa (MEY), Sản lượng tối đa hóa độ đàn hồi (RMY), Phân tích độ nhạy.

\vspace{1cm}
\hrule
\vspace{0.8cm}

\begin{center}
  {\Large \bfseries ABSTRACT}
\end{center}
\vspace{0.5cm}

Sustainable use of biological resources is very important as over-exploitation in the long run may lead to stock depletion, which in turn may threaten biodiversity. The Chesapeake Bay is an extremely complex ecosystem, and sustainable harvesting of its fisheries is essential both for the ecosystem's biodiversity and economic prosperity of the area. In this report, based on the research by Iordanka N. Panayotova, John Herrmann, and Nathan Kolling (published in \textit{Ecological Modelling}, 2023), we investigate an ecosystem-based mathematical modeling framework to study the population dynamics with harvesting of two key fish species in the Chesapeake Bay: the Atlantic Menhaden (\textit{Brevoortia tyrannus}) as prey and the Striped Bass (\textit{Morone saxatilis}) as predator.

The study fits the generalized Lotka--Volterra model to actual time series abundance data of the two species obtained from fisheries in the Bay from 1995 to 2017. Conditions for the existence of bioeconomic equilibria are derived, and the stability and resilience of the biological system are comprehensively investigated. Three fundamental harvesting policies are analyzed and compared: Maximum Sustainable Yield (MSY), Maximum Economic Yield (MEY), and Resilience Maximizing Yield (RMY), highlighting their long-term effects on fisheries sustainability. Furthermore, local sensitivity analysis using the direct differential method identifies the key biological parameters governing the population dynamics. This mathematical framework provides actionable insights for policymakers seeking to balance economic returns with ecosystem conservation in multi-species marine fisheries.

\vspace{0.4cm}
\noindent \textbf{Keywords:} Mathematical modeling, Differential equations, Predator--prey dynamics, Chesapeake Bay fisheries, Bioeconomics, Maximum sustainable yield, Maximum economic yield, Resilience maximizing yield, Sensitivity analysis.
